\subsection{Proof of the main theorem}\label{sec:main-theorem}

\begin{theorem}[main theorem]\label{thm:main-theorem}
  The functor
  \[
    \mathcal{D}:\mathbf{Int}\to\mathbf{DInt}
  \]
  is left adjoint to the inclusion functor
  \[
    \iota_{D}:\mathbf{DInt}\to\mathbf{Int} .
  \]
\end{theorem}
\begin{proof}
  We argue for the functor $\mathcal{D}$ defined in
  Lemma~\ref{lem:daniell-ext-is-functor-2}. Let $(X,L,I)$ be an arbitrary
  integration system and $(X,L^1,I^1)$ its Daniell extension.
  It suffices to show that the morphism of integration systems
  \[
    i_{(X,L,I)}:(X,L,I)\to(X,L^1,I^1)
  \]
  induced by the identity map $\operatorname{id}_X$ is a universal morphism from
  $(X,L,I)$ to $\iota_{D}$.

  Let $(Y,M,J)$ be a Daniell closed integration system on a set $Y$ and
  let
  \[
    F:(X,L,I)\to\iota_{D}(Y,M,J)
  \]
  be a morphism of integration systems. By
  Lemma~\ref{lem:daniell-ext-is-functor-1}, the map of base spaces
  $\underline{F}:Y\to X$ induces a morphism of integration systems
  \[
    F^1:(X,L^1,I^1)\to(Y,M^1,J^1) .
  \]

  Since $(Y,M,J)$ is Daniell closed we have $M^1=M$. Moreover
  $J^1(g)=J(g)$ for every $g\in M$ by
  Lemma~\ref{lem:daniell-ext-contains-source}. Hence
  \[
    (Y,M^1,J^1)=(Y,M,J) .
  \]
  Under this identification we set
  \[
    F^\sharp:=F^1:(X,L^1,I^1)\to(Y,M,J) .
  \]

  The inclusion functor does not change morphisms, and the map of base spaces
  underlying $i_{(X,L,I)}$ is $\operatorname{id}_X$, so the map of base spaces
  underlying the composite $\iota_{D}(F^\sharp)\circ i_{(X,L,I)}$ is
  \[
    \underline{i_{(X,L,I)}}\circ\underline{F^\sharp}
    =\operatorname{id}_X\circ\underline{F}
    =\underline{F} .
  \]
  A morphism of integration systems is precisely its map of base spaces, so
  \[
    F=\iota_{D}(F^\sharp)\circ i_{(X,L,I)} .
  \]

  Next, suppose a morphism
  \[
    H:(X,L^1,I^1)\to(Y,M,J)
  \]
  satisfies
  \[
    F=\iota_{D}(H)\circ i_{(X,L,I)} .
  \]
  Comparing the maps of base spaces on both sides,
  \[
    \underline{F}
    =\underline{i_{(X,L,I)}}\circ\underline{H}
    =\operatorname{id}_X\circ\underline{H}
    =\underline{H} .
  \]
  On the other hand $\underline{F^\sharp}=\underline{F}$, so $H=F^\sharp$.
  Hence $F^\sharp$ is the unique morphism satisfying the factorization above,
  and $i_{(X,L,I)}$ is a universal morphism from $(X,L,I)$ to $\iota_{D}$.

  Since $(X,L,I)$ was arbitrary,
  Theorem~\ref{thm:left-adjunction-existence-from-universality} yields a left
  adjoint of $\iota_{D}$, and its action on objects is the Daniell extension.
  Moreover, for every morphism of integration systems
  \[
    T:(X,L,I)\to(X',L',I')
  \]
  the morphism $\mathcal{D}(T)=T^1$ satisfies
  \[
    \iota_{D}(\mathcal{D}(T))\circ i_{(X,L,I)}
    =i_{(X',L',I')}\circ T ,
  \]
  since the maps of base spaces on both sides are $\underline{T}$.
  By the uniqueness coming from the universal property, the action on morphisms
  of the left adjoint constructed in that theorem agrees with the action of
  $\mathcal{D}$ on morphisms. Therefore $\mathcal{D}$ is left adjoint to
  $\iota_{D}$.
\end{proof}



\begin{remark}[on the content of the main theorem]\label{rem:main-theorem-not-tautology}
  The skeleton of the proof of the main theorem is formal.
  The pair consisting of $T:=\iota_{D}\circ\mathcal{D}$ and the morphisms $i_A$
  is an idempotent pointed endofunctor
  (one has $T(i_A)=i_{TA}=\operatorname{id}$), and it is a standard fact of
  category theory that the full subcategory of fixed points of an idempotent
  pointed endofunctor is reflective.
  Indeed, all that the proof above uses is the following two points:
  \begin{itemize}
    \item idempotence $(L^1)^1=L^1$
      (Lemma~\ref{lem:daniell-ext-is-complete});
    \item that pullback preserves integrability and the value of the integral
      (Lemma~\ref{lem:daniell-ext-is-functor-1}).
  \end{itemize}
  The former is stated as a corollary in \cite{Yamagami} and is known.

  \textbf{Hence, as long as $\mathbf{DInt}$ is defined as ``the category of
  integration systems satisfying $L^1=L$'', the main theorem says no more than
  that the subcategory of fixed points of $\mathcal{D}$ has $\mathcal{D}$ as its
  reflector, and carries no new content.}

  Let us say in advance that the following line of argument does not work: that
  the main theorem has content because
  Theorem~\ref{thm:complete-intrinsic} makes it possible to write Daniell
  closedness ``without mentioning $L^1$''.
  That is invalid. The condition $L^1=L$, when unfolded, also reads ``every
  function admitting a series representation by elements of $L$ belongs to
  $L$'', which uses nothing but $L$, $I$ and pointwise conditions
  (Remark~\ref{rem:not-about-mentioning-L1}).
  Being a reflective subcategory, too, follows formally under either
  description.
  If the rule of closure is preserved under pullback then the subcategory of
  closed objects is reflective, and the proof of
  Theorem~\ref{thm:mint-inclusion-has-left-adjoint} (the reflection into
  $L^{\mathrm{M}}$) has exactly that shape.

  The significance of Theorem~\ref{thm:complete-intrinsic} lies elsewhere.
  It is that \textbf{one closure operation, the Daniell extension, splits into
  two independent closure operations}:
  \[
    L^1=L
    \quad\Longleftrightarrow\quad
    \text{closure under the monotone convergence theorem}
    \quad\textbf{and}\quad
    \text{closure under modification on null sets} .
  \]
  Both conditions on the right occur independently in measure theory.
  The monotone convergence theorem corresponds directly, and closure under
  modification on null sets corresponds to completeness of the measure.
  That neither can be dropped was checked in
  Remark~\ref{rem:two-conditions-independent}.

  Under this decomposition, the main theorem asserts the following.

  \begin{quote}
    The full subcategory of $\mathbf{Int}$ cut out by the two conditions of the
    monotone convergence theorem and modification on null sets is reflective,
    and its reflector is given by the Daniell extension via series
    representations, a construction that can be written down
    \emph{explicitly}.
  \end{quote}

  \noindent
  The content thus moves away from the proof of the adjunction and over to the
  fact that the subcategory cut out by the two conditions coincides with the
  whole set of fixed points of $\mathcal{D}$
  (Theorem~\ref{thm:complete-intrinsic}).
  The two places in this paper with technical content are this one and
  Theorem~\ref{thm:LM-equals-Lsigma}; the proof of the adjunction itself is
  formal.

  This account concerns $\mathbf{Int}$ in
  Section~\ref{chap:functoriality} and says nothing about the category of
  operator morphisms in Section~\ref{chap:operator}.
\end{remark}
